\SourceLeaf{164}
machen: Deckung des Deficits kann und muss ganz allein durch zweckmässige Abänderungen in den bisherigen Einrichtungen bewirkt werden.

Es ist von selbst klar, dass schwache Mittel auch nur schwache Wirkungen hervorbringen können. Zu solchen wäre zu rechnen: die Verwandlung der bisher freiwilligen Theilnahme in eine gezwungene fur alle künftig ernannte Professoren, und die Aufhebung der Freiheit, zu jeder Zeit wieder auszutreten. Dass die Wirksamkeit einer solchen Maassregel nur eine sehr geringe sein könne, erhellet aus dem Umstande, dass, nach einem 50jährigen Durchschnitt, die mittlere Anzahl der nicht beitragenden obwohl zur Theilnahme berechtigten Professoren nur \(=4{,}4\) gewesen ist, also der Betrag der dadurch der Kasse jährlich entgehenden Einnahme, nach bisheriger Beitragshöhe \(=44\) Thaler, was mithin nach Verschiedenheit des Zinsfusses einem eisernen Kapital von 1100 oder von 1257 Rthl. gleichkommt. Allein diess wird, wenn auch nicht ganz, doch grossentheils durch die Strafgelder aufgewogen, welche nach der bestehenden sehr zweckmässigen Einrichtung bei verspäteten Beitritten zu erlegen und zum Theil sehr beträchtlich gewesen sind, wie aus folgenden Proben zu ersehen ist: Es haben doppelt nachgetragen
\begin{center}
L. für 12 Jahre,\quad R. für 13 Jahre,\quad L. für 19 Jahre,\quad S. für 25 Jahre,\quad \mbox{C. für 30 Jahre.}
\end{center}
\noindent\hspace*{\parindent}Die Commission ist daher der Meinung, dass ein so geringer und zweifelhafter Vortheil, wie aus der Verwandlung des Instituts in eine Zwangsanstalt für die Kasse hervorgehen könnte, gegen die Zerstörung des bisherigen liberalen Charakters dieser Stiftung nicht in Betracht kommen dürfe.

Ein paar andere Mittel von gleichfalls nur schwacher oder unsicherer Wirksamkeit werden am Schluss dieses Berichts erwähnt werden.

Als wirklich kräftige Mittel können demnach nur betrachtet werden:

1) Erhöhung der Beiträge.

2) Herabsetzung der Pensionen.

3) Verbindung beider Mittel.

Die in der Denkschrift aufgestellte Bilanzrechnung gewährt die Möglichkeit, genau anzugeben, in welchem Maasse diese Mittel in Anwendung gebracht werden müssen, wenn der Zweck erreicht werden, d. i. die Bilanz um 17520 Rthl. gebessert erscheinen soll. Details darüber würden hier nicht an ihrem Platze sein; das Endresultat aber ist:

1) Wenn das Deficit bloss durch Erhöhung der Beiträge gedeckt werden soll, so müssen diese allgemein, d. i. für \textit{alle} jetzigen und künftigen Mitglieder, auf \(4\tfrac{1}{4}\) Louisd’or erhöhet werden.

2) Soll die Herabsetzung der Pensionen allein die Deckung bewirken, so müssen dieselben auf 223 Rthl. reducirt werden (wobei der der Professorin G. auf 200 Rthl. bestehen bliebe).

3) Bei einer Vertheilung der Last auf Beitragende und Pensionirte käme es darauf an, welches Verhältniss der Theilung man wählte: sollten z. B. erstere \(\tfrac{2}{3}\), letztere \(\tfrac{1}{3}\) übernehmen, so würden die Beiträge \(3\tfrac{1}{2}\) Louisd’or, die Pensionen 241 Rthl. betragen müssen. Sollen die Beiträge nur auf 3 Louisd’or erhöhet werden, so können, wenn das Deficit wirklich gedeckt sein soll, nur 235 Rthl. Pension verabreicht werden.

Bei allen diesen Rechnungen ist vorausgesetzt, dass die gewählten Änderungen von Michaelis 1845 an in Wirksamkeit treten, und dass im zweiten oder dritten Falle die Pensionsherabsetzungen sämmtliche Witwen, die gegenwärtigen wie die künftigen treffen. Sollten, ohne alle Beitragserhöhung und ohne Herabsetzung der Pension für die gegenwärtigen Witwen, die künftigen Witwen die Last allein tragen, so könnte für diese nur die Pension zu \(213\tfrac{1}{4}\) Rthl. gewährt werden.
